% !TEX TS-program = xelatex
% !TEX encoding = UTF-8 Unicode
% !Mode:: "TeX:UTF-8"

\documentclass{resume}
\usepackage{zh_CN-Adobefonts_external} % Simplified Chinese Support using external fonts (./fonts/zh_CN-Adobe/)
%\usepackage{zh_CN-Adobefonts_internal} % Simplified Chinese Support using system fonts
\usepackage{linespacing_fix} % disable extra space before next section
\usepackage{cite}

\usepackage{pifont} % 圈数字符号

\begin{document}
\pagenumbering{gobble} % suppress displaying page number

\name{孙铎}

% {E-mail}{mobilephone}{homepage}
% be careful of _ in emaill address
\contactInfo{中共党员}{15818521182}{1075079562@qq.com}{GitHub @d-sun-cs}
% {E-mail}{mobilephone}
% keep the last empty braces!
%\contactInfo{xxx@yuanbin.me}{(+86) 131-221-87xxx}{}
 
% \section{个人总结}
% 本人在校成绩优秀、乐观向上，工作负责、自我驱动力强、热爱尝试新事物，认同开放、连接、共享的Web在未来的不可替代性。在校期间长期从事可视分析(Web的2D/3D时空可视化)相关研究，对Web技术发展趋势及前端工程化解决方案有浓厚兴趣。\textbf{现任职于 BAT 集团。}

% \section{\faGraduationCap\ 教育背景}
\section{教育背景}
\datedsubsection{\textbf{哈尔滨工业大学（深圳）}| 计算机科学与技术 | 学硕在读}{2024.9 - 至今}
\begin{itemize}
  \item \textbf{绩点3.657/4} | 研究方向：高性能计算、并行计算 | 课程分数/100：矩阵分析/99、高级算法/92.4、并行处理/91.5
\end{itemize}
\datedsubsection{\textbf{哈尔滨工业大学（深圳）}| 计算机科学与技术 | 学士}{2020.9 - 2024.6}
\begin{itemize}
  \item \textbf{学分绩94.7/100} | \textbf{排名10/527(前2\%)} | \textbf{国家奖学金（两次）}| 国家励志奖学金 | 华为智能基座奖学金 | 哈尔滨工业大学优秀团员/学生标兵 | 哈尔滨工业大学优秀毕业生 | 课程分数/100：操作系统/96、计算机网络/97
\end{itemize}

\section{技术能力}
\begin{itemize}[parsep=0.2ex]
  \item \textbf{编程语言}: 熟悉C/C++、CUDA，了解Python、Shell、Java、JS/HTML/CSS、Go、SQL
  \item \textbf{工程构建}: 熟悉CMake、Git、Linux、Docker，了解Redis/MySQL、Webpack、Gradle、Maven
  \item \textbf{开发框架}: 熟悉cuBLAS、cuSparse了解CUTLASS、Spring（Boot | Security）、MyBatis、Vue.js、Chrome Extension
\end{itemize}


\section{实习经历}
\datedsubsection{\textbf{天王星量化深圳研究院}}{2024.06 - 2024.09}
\textit{推理框架优化项目 | C++语言 | 高性能计算开发实习生}
\begin{itemize}
  \item 背景：公司专用量化模型的推理框架有待进一步优化，并需要完善文档以支撑后续迭代。
  \item 任务：为推理框架撰写技术文档，学习开源推理框架（ONNX Runtime），进一步探索计算图的推理加速方案。
  \item 工作：\ding{172}模块化撰写项目文档，了解框架代码基本情况；\ding{173}提出“懒执行”优化策略，保存输入节点最大子树以实现增量执行；\ding{174}设计线程池并发执行算子，为张量绑定条件变量，解决DAG同步问题。
  \item 结果：测试场景下，“懒执行”提速约20\%，线程池并发推理在4-8核CPU下提速约120\%。
\end{itemize}

\section{项目经历}
\datedsubsection{\textbf{移位异或纠删码的Ceph集成与优化}}{2025.02 - 2025.08}
\textit{硕士科研项目/毕业设计 | C++语言}
\begin{itemize}
  \item 背景：导师自研的移位异或纠删码理论性能高于传统纠删码，实际性能可在存储系统中进一步验证。
  \item 任务：将移位异或纠删码集成至开源分布式存储系统Ceph，做出一定的适配与优化，并测试其性能。
  \item 工作：\ding{172}分析Ceph纠删码结构，实现其接口以集成新算法；\ding{173}用Redis存储移位多出的校验数据（Ceph只能存对齐部分）；\ding{174}实现错位窗口与相邻合并优化内存访问，提出“纠删映射”高效泛化解码。
  \item 结果：成功通过Ceph内置纠删码功能测试；编解码相较于Ceph中SOTA算法ISA提速约30\%。
\end{itemize}
\datedsubsection{\textbf{面向新型GPU架构的稀疏矩阵乘法（SpMM）加速器}}{2023.11 - 2024.05}
\textit{本科科研项目/毕业设计 | CUDA C++语言}
\begin{itemize}
  \item 背景：内存访问是SpMM的性能瓶颈，Ampere架构新增共享内存相关特性，可用于缓解内存瓶颈
  \item 任务：利用从全局内存到共享内存异步数据拷贝的新特性，通过流水线数据预取优化SpMM算法
  \item 工作：用ncu分析性能瓶颈，用新特性改进现有研究算法，实现适配新特性的复用稠密矩阵算法
  \item 结果：应用新特性的算法在小矩阵“隐藏延时”不足的情况下性能较官方库更好，项目研究仍在继续
\end{itemize}
\datedsubsection{\textbf{基于SPDK实现Linux用户态IO多路径软件}}{2023.03 - 2023.08}
\textit{全国大学生计算机系统能力大赛与华为公司联合赛题 | 获评国家级一等奖 | 参赛队长 | C语言}
\begin{itemize}
  \item 背景：SPDK是新型用户态高性能存储驱动，但是其绕过了Linux内核，无法使用内核多路径功能
  \item 任务：模仿内核多路径功能，基于SPDK开发用户态多路径软件，用于高效稳定的虚拟化存储网络
  \item 工作：用动态链接库添加用户态多路径，实现故障处理和负载均衡等功能，编写项目文档
  \item 结果：存储网络中，添加多路径的SPDK在稳定性和性能方面优于单路径和内核多路径。
\end{itemize}
\datedsubsection{\textbf{古代玻璃制品的成分分析与鉴别}}{2022.09 - 2022.09}
\textit{全国大学生数学建模竞赛赛题 | 获评国家级二等奖 | 参赛队长 | Python语言}
\begin{itemize}
  \item 背景：古代玻璃文物有高钾和铅钡两类，受风化影响其化学成分有所变化，理化性质呈现出一定规律
  \item 任务：根据题目给出的化学成分数据，分析玻璃文物各种“理化特征”之间的统计规律与关联关系
  \item 工作：用统计图、卡方检验、决策树建立玻璃分类模型，用K-means聚类进行亚类划分，书写论文
  \item 结果：正确预测成分含量和样本分类，合理划分出亚类，找到成分关联关系，完成结果的检验分析。 
\end{itemize}

\section{课余活动}
\begin{itemize}[parsep=0.2ex]
  \item 学院第四届学生会主席：三位主席之一，举办生产力工具分享等活动，荣获深圳市优秀学生会称号
  \item 校艺术团电声部乐队键盘手：多次参加晚会表演，担任学院合唱团钢琴伴奏，取得校合唱比赛冠军
\end{itemize}

\end{document}
